% from ega2-indexes-and-contents-fr.tex
\oldpage[II]{207}
\markboth{A. GROTHENDIECK --- Chap. II}{A. GROTHENDIECK --- Chap. II}
\phantomsection
\addcontentsline{toc}{section}{Index des notations}
\section*{Index des notations}
\thispagestyle{plain}

\begingroup
\small
\setlength{\parindent}{0pt}
\setlength{\parskip}{2pt}
\newcommand{\egaShProj}{\mathcal{P}\!\operatorname{roj}}

$\mathcal{A}(X),\mathcal{A}(\mathcal{F}),\mathcal{A}(\mathcal{B})$
($X$ $S$-préschéma, $\mathcal{F}$ $\mathcal{O}_X$-Module,
$\mathcal{B}$ $\mathcal{O}_X$-Algèbre) : \hyperref[II.1.1.1-fr]{1.1.1}.

$\mathcal{A}(h)$ ($h$ $S$-morphisme) : \hyperref[II.1.1.2-fr]{1.1.2}.

$\Spec(\mathcal{B})$ ($\mathcal{B}$ $\mathcal{O}_S$-Algèbre) : \hyperref[II.1.3.1-fr]{1.3.1}.

$\widetilde{\mathcal{M}}$ ($\mathcal{M}$ $\mathcal{B}$-Module,
$\mathcal{B}$ $\mathcal{O}_S$-Algèbre) : \hyperref[II.1.4.3-fr]{1.4.3}.

$\mathbf{T}(E),\mathbf{S}(E),\mathbf{S}_A(E)$ ($E$ $A$-module) : \hyperref[II.1.7.1-fr]{1.7.1}.

$\mathbf{S}(\mathcal{E}),\mathbf{S}_{\mathcal{A}}(\mathcal{E})$
($\mathcal{E}$ $\mathcal{A}$-Module) : \hyperref[II.1.7.4-fr]{1.7.4}.

$\mathbf{V}(\mathcal{E})$ ($\mathcal{E}$ $\mathcal{O}_S$-Module quasi-cohérent) : \hyperref[II.1.7.8-fr]{1.7.8}.

$S_n,S_+,M_n,S^{(d)},M^{(d,k)},M^{(d)}$
($S$ anneau gradué, $M$ $S$-module gradué) : \hyperref[II.2.1.1-fr]{2.1.1}.

$M(n),S(n)$ ($S$ anneau gradué, $M$ $S$-module gradué) : \hyperref[II.2.1.1-fr]{2.1.1}.

$\Hom_S(M,N)$ ($S$ anneau gradué, $M,N$ $S$-modules gradués) : \hyperref[II.2.1.2-fr]{2.1.2}.

$\mathfrak N_+,\mathfrak r_+(\mathfrak J)$
($\mathfrak J$ idéal de $S_+$) : \hyperref[II.2.1.10-fr]{2.1.10}.

$S_{(f)},M_{(f)}$ ($f$ élément homogène de $S_+$) : \hyperref[II.2.2.1-fr]{2.2.1}.

$S_{(T)},M_{(T)},S_{(\mathfrak p)},M_{(\mathfrak p)}$
($T$ partie multiplicative de $S_+$ formée d'éléments homogènes,
$\mathfrak p$ idéal premier gradué de $S_+$) : \hyperref[II.2.2.7-fr]{2.2.7}.

$\Proj(S)$ ($S$ anneau gradué) : \hyperref[II.2.3.1-fr]{2.3.1}.

$V_+(E)$ ($E$ partie de $S_+$) : \hyperref[II.2.3.2-fr]{2.3.2}.

$D_+(f)$ ($f$ élément de $S$) : \hyperref[II.2.3.3-fr]{2.3.3}.

$\mathfrak j_+(Y)$ ($Y$ partie de $\Proj(S)$) : \hyperref[II.2.3.10-fr]{2.3.10}.

$\widetilde M$ ($M$ $S$-module gradué) : \hyperref[II.2.5.2-fr]{2.5.2}.

$\widetilde u$ ($u$ homomorphisme de degré zéro) : \hyperref[II.2.5.4-fr]{2.5.4}.

$\mathcal{O}_X(n),\mathcal{F}(n)$
($X=\Proj(S)$, $\mathcal{F}$ $\mathcal{O}_X$-Module) : \hyperref[II.2.5.10-fr]{2.5.10}.

$\lambda$ (homomorphisme canonique) : \hyperref[II.2.5.11-fr]{2.5.11}.

$\mu$ (homomorphisme canonique) : \hyperref[II.2.5.12-fr]{2.5.12}.

$X^{(d)}$ ($X=\Proj(S)$, $d$ entier $>0$) : \hyperref[II.2.5.16-fr]{2.5.16}.

$\Gamma_*(\mathcal{F})$ ($\mathcal{F}$ $\mathcal{O}_X$-Module,
$X=\Proj(S)$) : \hyperref[II.2.6.1-fr]{2.6.1}.

$\alpha_n,\alpha,\alpha_{\mathcal{M}}$ (homomorphismes canoniques) : \hyperref[II.2.6.2-fr]{2.6.2}.

$\beta,\beta_{\mathcal{F}}$ (homomorphismes canoniques) : \hyperref[II.2.6.4-fr]{2.6.4}.

$G(\vphi),{}^a\vphi$ (par abus de langage)
($\vphi$ homomorphisme d'anneaux gradués) : \hyperref[II.2.8.1-fr]{2.8.1}.

$\Proj(\vphi)$ ($\vphi$ homomorphisme d'anneaux gradués) : \hyperref[II.2.8.2-fr]{2.8.2}.

$\mathcal{S}^{(d)},\mathcal{M}^{(d,k)},\mathcal{M}^{(d)},\mathcal{M}(n)$
($\mathcal{S}$ $\mathcal{O}_Y$-Algèbre graduée,
$\mathcal{M}$ $\mathcal{S}$-Module gradué) : \hyperref[II.3.1.1-fr]{3.1.1}.

$\Proj(\mathcal{S})$ ($\mathcal{S}$ $\mathcal{O}_Y$-Algèbre graduée quasi-cohérente) : \hyperref[II.3.1.3-fr]{3.1.3}.

$X_f$ ($X=\Proj(\mathcal{S})$, $f$ section de $\mathcal{S}_d$ au-dessus de $Y$) : \hyperref[II.3.1.4-fr]{3.1.4}.

$\widetilde{\mathcal{M}}$ ($\mathcal{M}$ $\mathcal{S}$-Module gradué) : \hyperref[II.3.2.2-fr]{3.2.2}.

$\mathcal{O}_X(n),\mathcal{F}(n)$
($X=\Proj(\mathcal{S})$, $\mathcal{F}$ $\mathcal{O}_X$-Module) : \hyperref[II.3.2.5-fr]{3.2.5}.

$\lambda,\mu$ (homomorphismes canoniques) : \hyperref[II.3.2.6-fr]{3.2.6}.

$\Gamma_*(\mathcal{F})$ ($\mathcal{F}$ $\mathcal{O}_X$-Module,
$X=\Proj(\mathcal{S})$) : \hyperref[II.3.3.1-fr]{3.3.1}.

$\alpha_n,\alpha,\alpha_{\mathcal{M}}$ (homomorphismes canoniques) : \hyperref[II.3.3.2-fr]{3.3.2}.

